\chapter{DX-012 --- Document Render-Path Naming Convention}\label{record-344}

\noindent{\footnotesize\color{BookMuted}\path{docs/method/retro/DX_012-render-path-naming-convention.md}\par}\smallskip

\begin{RecordMeta}
title: DX-012 — Document Render-Path Naming Convention
lane: retro
legend: DX

\end{RecordMeta}

\section{Disposition}\label{record-344-disposition}

Documented the render-path naming rule explicitly in the \texttt{@\allowbreak{}f\allowbreak{}l\allowbreak{}y\allowbreak{}i\allowbreak{}n\allowbreak{}g\allowbreak{}r\allowbreak{}o\allowbreak{}b\allowbreak{}o\allowbreak{}t\allowbreak{}s\allowbreak{}/\allowbreak{}b\allowbreak{}i\allowbreak{}j\allowbreak{}o\allowbreak{}u} README and GUIDE: the base family name is the public component family, and a \texttt{*\allowbreak{}S\allowbreak{}u\allowbreak{}r\allowbreak{}f\allowbreak{}a\allowbreak{}c\allowbreak{}e\allowbreak{}(\allowbreak{})} companion means the same family on the composable surface path. The broader doctrine was already present across ADVANCED\_GUIDE, migration docs, and design-system material, so this closes the remaining front-door POLA gap the note described.

\section{Original Proposal}\label{record-344-original-proposal}

Legend: \href{https://github.com/flyingrobots/bijou/blob/1853ad902f6bf219326409e7f743a5a1d6517dbc/docs/method/legends/DX-developer-experience.md}{DX --- Developer Experience}

\section{Idea}\label{record-344-idea}

Document the implicit naming convention for string-output vs Surface-output component variants:

\begin{itemize}
\tightlist
\item
  \texttt{b\allowbreak{}o\allowbreak{}x\allowbreak{}(\allowbreak{}c\allowbreak{}o\allowbreak{}n\allowbreak{}t\allowbreak{}e\allowbreak{}n\allowbreak{}t\allowbreak{})} → string (pipe-friendly)
\item
  \texttt{b\allowbreak{}o\allowbreak{}x\allowbreak{}S\allowbreak{}u\allowbreak{}r\allowbreak{}f\allowbreak{}a\allowbreak{}c\allowbreak{}e\allowbreak{}(\allowbreak{}c\allowbreak{}o\allowbreak{}n\allowbreak{}t\allowbreak{}e\allowbreak{}n\allowbreak{}t\allowbreak{})} → Surface (composable in TUI apps)
\item
  \texttt{b\allowbreak{}a\allowbreak{}d\allowbreak{}g\allowbreak{}e\allowbreak{}(\allowbreak{}c\allowbreak{}o\allowbreak{}n\allowbreak{}t\allowbreak{}e\allowbreak{}n\allowbreak{}t\allowbreak{})} → Surface only (no string variant)
\item
  \texttt{t\allowbreak{}a\allowbreak{}b\allowbreak{}l\allowbreak{}e\allowbreak{}(\allowbreak{}d\allowbreak{}a\allowbreak{}t\allowbreak{}a\allowbreak{})} → string (mixed-mode, pipe-aware)
\end{itemize}

The pattern exists but is never stated explicitly. A developer browsing the API doesn't know why some components have a \texttt{-\allowbreak{}S\allowbreak{}u\allowbreak{}r\allowbreak{}f\allowbreak{}a\allowbreak{}c\allowbreak{}e} variant and others don't, or why \texttt{b\allowbreak{}a\allowbreak{}d\allowbreak{}g\allowbreak{}e\allowbreak{}(\allowbreak{})} returns a Surface directly while \texttt{b\allowbreak{}o\allowbreak{}x\allowbreak{}(\allowbreak{})} returns a string.

\section{Why}\label{record-344-why}

This came up in the DX audit as a minor POLA concern. The convention is sound (simpler output-only components are string-first; richer interactive components are Surface-native), but it reads as inconsistency when you're scanning exports without context.

A one-paragraph note in the bijou README's component section would resolve the confusion.

